\scanpage{24}
\begingroup\csname tagsleft@true\endcsname
\[
\deleted{\Delta_A(t)=\Delta_{\lvert A\rvert}(t)=\Delta_{\varphi_{\lvert A\rvert}}(t)
=\exp\int_0^t\log\varphi_{\lvert A\rvert}(s)\,ds}
\tag*{\deleted{(5)}}
\]
\[
\deleted{\Delta(A)=\deleted{\unsure{\int_0^{T(1)}}}\Delta_A(\omega)
\qquad(\text{si }\{\omega=T(1)\}<+\infty)}
\tag*{\deleted{(6)}}
\]
\endgroup

\deleted{\unsure{Or|$\varphi$} Par suite, $\Delta_A(t)$ est une fonction \deleted{\unsure{finie}} continue de $t$ à valeurs $0\leq\Delta_A(t)<\infty$. $\Delta(A)$ est \deleted{\unsure{tjrs}} finie si $\unsure{\mathrm{m.}}T(1)<\infty$ (i.e. si $T$ est une trace finie), \unsure{sinon} \unsure{Si|si} la trace est \unsure{finie|bornée}, on a $1\leq\Delta(A)<+\infty$.}

% slanted
\margin{\unsure{Définition} de \deleted{\unsure{de la forme}} $\varphi_f=\unsure{1+f|1+f^+}$, $f\in\tilde b$ $(\unsure{\bar b|\tilde b}=C+\underline b)$. Définition de $\unsure{W|W^*}(A)$. les ``$A$'' \unsure{$\in$|de} l'idéal \deleted{\unsure{de finie}} $\unsure{\underline a_0|a_0}$ pour \unsure{$\equiv$}} Supposons que $\underline C$ soit une sous-algèbre commutative \unsure{auto-adjointe} maximale de $\underline b$, $\underline C=\mathcal C_0(M)$, $m$ la mesure sur $M$ définie par la \unsure{restriction de|restriction à} $T$. Si $\unsure{t|s}\in M$ est un point de $M$ de mesure non nulle, on voit \unsure{bien|par} \unsure{(en vertu de|(l'étude de} la maximalité de $\underline C$) que \inserted{$E=\unsure{U_{\varphi_{\{t\}}}|\!U\varphi_{\{t\}}}$} \deleted{\unsure{$E_\varphi$}} est un projecteur minimal de \deleted{\unsure{$\mathcal A$}} \unsure{$\underline b$} \unsure{Réciproquement} \unsure{soit|si} $E$ est un projecteur minimal dans \inserted{\unsure{$\mathcal A$}} \deleted{\unsure{$\underline E$}, $\underline C$ \inserted{$=\mathcal C_0(M)$} une algèbre \unsure{autoadjointe} maximale de $\underline b$, contenant $E$, alors \deleted{il est} \deleted{\unsure{immédiat}} \unsure{identique} à la fonction \deleted{$E$ est la fonction} caractéristique d'un \deleted{point} \deleted{ensemble à la fois} \deleted{ouvert et \unsure{fermé}} \unsure{ensemble} ouvert et fermé de $M$, \deleted{\unsure{vide} formé} \unsure{évident|réduit} à un point donc d'une \inserted{donc} mesure $>0$. \unsure{Ceci|On} dit \deleted{\unsure{qu'un} \unsure{l'opérateur}} \inserted{que la} \inserted{trace} $T$ est \deleted{\unsure{continue}} non atomique si il n'existe pas de projecteur $>0$ minimal dans $\underline{\unsure{C|b}}$, i.e. si \deleted{\unsure{on ne}} la mesure $\unsure{m|\mu}$ sur un espace $M$. i.e. de à \unsure{une} \unsure{$\underline C$} \inserted{$=\mathcal C_0(M)$} \unsure{associé est} non atomique \unsure{tout} \unsure{$C^*$} \unsure{fermé} de $\mathcal A$.} \inserted{Si $E$ \uline{est} un} projecteur $>0$ \deleted{\unsure{nul}} minimal, \deleted{\unsure{dans un}} on voit facilement qu'il existe un tel
% archive sheet number: 25

